%%%%%%%%%%%%%%%%%%%%%%%%%%%%%%%%%%%%%%%%%%%%%%%%%%%%%%%%%%%%
% 8.2 DATA COLLECTION AND SAMPLING
% The Composition of Advanced Mathematics
%%%%%%%%%%%%%%%%%%%%%%%%%%%%%%%%%%%%%%%%%%%%%%%%%%%%%%%%%%%%

\section{Data Collection and Sampling}
\label{sec:data-collection-sampling}

\prerequisite{
Descriptive statistics, populations and samples, categorical and
quantitative variables, basic probability, conditional probability,
random variables, expectation, variance, and elementary set notation.
}

\learningobjectives{
By the end of this section, the reader should be able to:
\begin{itemize}
    \item distinguish observational studies from experiments;
    \item identify populations, target populations, sampling frames, and
          samples;
    \item distinguish probability and nonprobability sampling;
    \item construct simple random samples;
    \item distinguish sampling with and without replacement;
    \item understand stratified sampling;
    \item understand cluster and multistage sampling;
    \item recognize systematic sampling;
    \item understand unequal-probability sampling conceptually;
    \item distinguish sampling error from nonsampling error;
    \item recognize undercoverage and selection bias;
    \item recognize voluntary-response bias;
    \item recognize nonresponse bias;
    \item recognize response and measurement bias;
    \item distinguish confounding from sampling bias;
    \item understand random assignment;
    \item distinguish random sampling from random assignment;
    \item identify explanatory and response variables;
    \item identify treatments, experimental units, and control groups;
    \item understand replication, randomization, and blocking;
    \item distinguish placebo, blinding, and double blinding;
    \item understand matched-pairs designs;
    \item distinguish cross-sectional, longitudinal, retrospective, and
          prospective studies;
    \item understand census data and administrative data;
    \item interpret sampling weights;
    \item understand finite-population corrections;
    \item distinguish design-based and model-based perspectives;
    \item identify ethical issues in data collection;
    \item evaluate whether a study design supports descriptive,
          population, or causal conclusions;
    \item connect data collection with sampling distributions,
          estimation, experimental design, and statistical inference.
\end{itemize}
}

%%%%%%%%%%%%%%%%%%%%%%%%%%%%%%%%%%%%%%%%%%%%%%%%%%%%%%%%%%%%
\subsection{Why Data Collection Matters}
\label{subsec:why-data-collection-matters}
%%%%%%%%%%%%%%%%%%%%%%%%%%%%%%%%%%%%%%%%%%%%%%%%%%%%%%%%%%%%

Statistical analysis cannot be separated from the process used to obtain
the data.

A sophisticated model applied to poorly collected data may produce
precise but misleading conclusions.

Thus,

\[
\boxed{
\text{good statistical inference begins with good data collection}.
}
\]

Before analyzing a dataset, one should ask:

\begin{itemize}
    \item Where did the data come from?
    \item Who or what was measured?
    \item How were observations selected?
    \item What variables were measured?
    \item Were treatments assigned?
    \item What populations can the results represent?
    \item What causal conclusions, if any, are justified?
\end{itemize}

%%%%%%%%%%%%%%%%%%%%%%%%%%%%%%%%%%%%%%%%%%%%%%%%%%%%%%%%%%%%
\subsection{Target Population}
\label{subsec:target-population}
%%%%%%%%%%%%%%%%%%%%%%%%%%%%%%%%%%%%%%%%%%%%%%%%%%%%%%%%%%%%

\begin{definitionbox}{Target Population}
The target population is the complete collection of units to which the
investigator ultimately wishes to generalize conclusions.
\end{definitionbox}

Examples include:

\[
\text{all registered voters in a state},
\]

\[
\text{all manufactured components from a production process},
\]

or

\[
\text{all students enrolled at a university}.
\]

The target population is defined by the research question, not merely by
the available data.

%%%%%%%%%%%%%%%%%%%%%%%%%%%%%%%%%%%%%%%%%%%%%%%%%%%%%%%%%%%%
\subsection{Study Population}
\label{subsec:study-population}
%%%%%%%%%%%%%%%%%%%%%%%%%%%%%%%%%%%%%%%%%%%%%%%%%%%%%%%%%%%%

The study population is the population that is actually accessible to
the investigator.

It may differ from the target population.

For example, the target population may be

\[
\text{all adults in a region},
\]

while the accessible study population may include only

\[
\text{adults with listed telephone numbers}.
\]

Differences between these populations can introduce bias.

%%%%%%%%%%%%%%%%%%%%%%%%%%%%%%%%%%%%%%%%%%%%%%%%%%%%%%%%%%%%
\subsection{Sampling Frame}
\label{subsec:sampling-frame}
%%%%%%%%%%%%%%%%%%%%%%%%%%%%%%%%%%%%%%%%%%%%%%%%%%%%%%%%%%%%

\begin{definitionbox}{Sampling Frame}
A sampling frame is the operational list or mechanism used to identify
units eligible for selection into a sample.
\end{definitionbox}

Examples include:

\[
\text{a voter-registration list},
\]

\[
\text{a school enrollment database},
\]

or

\[
\text{a list of households}.
\]

A sampling frame should approximate the target population as closely as
possible.

%%%%%%%%%%%%%%%%%%%%%%%%%%%%%%%%%%%%%%%%%%%%%%%%%%%%%%%%%%%%
\subsection{Sampling Unit}
\label{subsec:sampling-unit}
%%%%%%%%%%%%%%%%%%%%%%%%%%%%%%%%%%%%%%%%%%%%%%%%%%%%%%%%%%%%

A sampling unit is an entity selected during the sampling process.

Depending on the design, a sampling unit may be:

\[
\text{a person},
\]

\[
\text{a household},
\]

\[
\text{a classroom},
\]

\[
\text{a hospital},
\]

or

\[
\text{a geographic region}.
\]

The sampling unit need not always be the same as the final observational
unit.

%%%%%%%%%%%%%%%%%%%%%%%%%%%%%%%%%%%%%%%%%%%%%%%%%%%%%%%%%%%%
\subsection{Observational Unit}
\label{subsec:observational-unit}
%%%%%%%%%%%%%%%%%%%%%%%%%%%%%%%%%%%%%%%%%%%%%%%%%%%%%%%%%%%%

An observational unit is the entity on which measurements are ultimately
recorded.

For example, schools may be sampled first, but students inside those
schools may be the observational units.

This distinction becomes important in cluster and multistage sampling.

%%%%%%%%%%%%%%%%%%%%%%%%%%%%%%%%%%%%%%%%%%%%%%%%%%%%%%%%%%%%
\subsection{Census}
\label{subsec:census}
%%%%%%%%%%%%%%%%%%%%%%%%%%%%%%%%%%%%%%%%%%%%%%%%%%%%%%%%%%%%

\begin{definitionbox}{Census}
A census attempts to collect information from every member of the
population.
\end{definitionbox}

A census eliminates ordinary sampling variability arising from selecting
only part of the population.

However, it may still suffer from:

\[
\boxed{
\text{nonresponse},
}
\]

\[
\boxed{
\text{measurement error},
}
\]

\[
\boxed{
\text{coverage error},
}
\]

and

\[
\boxed{
\text{data-processing errors}.
}
\]

Thus a census is not automatically error-free.

%%%%%%%%%%%%%%%%%%%%%%%%%%%%%%%%%%%%%%%%%%%%%%%%%%%%%%%%%%%%
\subsection{Sample Survey}
\label{subsec:sample-survey}
%%%%%%%%%%%%%%%%%%%%%%%%%%%%%%%%%%%%%%%%%%%%%%%%%%%%%%%%%%%%

A sample survey collects information from only part of the target
population.

Sampling can reduce:

\[
\text{cost},
\]

\[
\text{time},
\]

and

\[
\text{administrative burden}.
\]

A carefully designed probability sample may provide more reliable
information than a poorly executed census.

%%%%%%%%%%%%%%%%%%%%%%%%%%%%%%%%%%%%%%%%%%%%%%%%%%%%%%%%%%%%
\subsection{Probability Sampling}
\label{subsec:probability-sampling}
%%%%%%%%%%%%%%%%%%%%%%%%%%%%%%%%%%%%%%%%%%%%%%%%%%%%%%%%%%%%

\begin{definitionbox}{Probability Sample}
A probability sample is selected using a random mechanism in which each
eligible population unit has a known, nonzero probability of selection,
or one that can be determined from the sampling design.
\end{definitionbox}

Probability sampling provides a mathematical basis for quantifying
sampling variability.

Important designs include:

\[
\boxed{
\text{simple random sampling},
}
\]

\[
\boxed{
\text{stratified sampling},
}
\]

\[
\boxed{
\text{cluster sampling},
}
\]

and

\[
\boxed{
\text{multistage sampling}.
}
\]

%%%%%%%%%%%%%%%%%%%%%%%%%%%%%%%%%%%%%%%%%%%%%%%%%%%%%%%%%%%%
\subsection{Nonprobability Sampling}
\label{subsec:nonprobability-sampling}
%%%%%%%%%%%%%%%%%%%%%%%%%%%%%%%%%%%%%%%%%%%%%%%%%%%%%%%%%%%%

Nonprobability sampling does not use a known random selection mechanism.

Examples include:

\[
\boxed{
\text{convenience samples},
}
\]

\[
\boxed{
\text{voluntary-response samples},
}
\]

\[
\boxed{
\text{quota samples},
}
\]

and

\[
\boxed{
\text{judgment samples}.
}
\]

Such samples can be useful for exploratory work but may not support
standard design-based population inference.

%%%%%%%%%%%%%%%%%%%%%%%%%%%%%%%%%%%%%%%%%%%%%%%%%%%%%%%%%%%%
\subsection{Convenience Sampling}
\label{subsec:convenience-sampling}
%%%%%%%%%%%%%%%%%%%%%%%%%%%%%%%%%%%%%%%%%%%%%%%%%%%%%%%%%%%%

A convenience sample selects units that are easiest to access.

Examples include:

\[
\text{surveying students in one nearby classroom},
\]

or

\[
\text{polling shoppers at one store entrance}.
\]

Convenience sampling is easy to conduct but can strongly overrepresent
certain types of individuals.

%%%%%%%%%%%%%%%%%%%%%%%%%%%%%%%%%%%%%%%%%%%%%%%%%%%%%%%%%%%%
\subsection{Voluntary-Response Sampling}
\label{subsec:voluntary-response-sampling}
%%%%%%%%%%%%%%%%%%%%%%%%%%%%%%%%%%%%%%%%%%%%%%%%%%%%%%%%%%%%

In a voluntary-response sample, individuals decide for themselves whether
to participate.

Examples include:

\[
\text{online polls},
\]

\[
\text{call-in surveys},
\]

and

\[
\text{optional customer feedback forms}.
\]

People with strong opinions may be especially likely to respond.

This creates voluntary-response bias.

%%%%%%%%%%%%%%%%%%%%%%%%%%%%%%%%%%%%%%%%%%%%%%%%%%%%%%%%%%%%
\subsection{Simple Random Sampling}
\label{subsec:simple-random-sampling}
%%%%%%%%%%%%%%%%%%%%%%%%%%%%%%%%%%%%%%%%%%%%%%%%%%%%%%%%%%%%

\begin{definitionbox}{Simple Random Sample}
A simple random sample of size \(n\), abbreviated SRS, from a finite
population of size \(N\) is a sample selected so that every subset of
\(n\) distinct population units has the same probability of selection.
\end{definitionbox}

The number of possible samples is

\[
\boxed{
\binom{N}{n}.
}
\]

Therefore each possible SRS has probability

\[
\boxed{
\frac1{\binom{N}{n}}.
}
\]

%%%%%%%%%%%%%%%%%%%%%%%%%%%%%%%%%%%%%%%%%%%%%%%%%%%%%%%%%%%%
\subsection{Worked Example: SRS Probability}
\label{subsec:worked-srs-probability}
%%%%%%%%%%%%%%%%%%%%%%%%%%%%%%%%%%%%%%%%%%%%%%%%%%%%%%%%%%%%

\begin{workedexamplebox}{Selecting Three Students}

Suppose a class contains

\[
N=10
\]

students.

A simple random sample of

\[
n=3
\]

students is selected.

The number of possible samples is

\[
\binom{10}{3}
=
120.
\]

Therefore any particular three-student sample has probability

\begin{answerbox}
\[
\boxed{
\frac1{120}.
}
\]
\end{answerbox}

\end{workedexamplebox}

%%%%%%%%%%%%%%%%%%%%%%%%%%%%%%%%%%%%%%%%%%%%%%%%%%%%%%%%%%%%
\subsection{Inclusion Probability Under SRS}
\label{subsec:srs-inclusion-probability}
%%%%%%%%%%%%%%%%%%%%%%%%%%%%%%%%%%%%%%%%%%%%%%%%%%%%%%%%%%%%

For an SRS of size \(n\) from a population of size \(N\), each individual
population unit has probability

\[
\boxed{
\frac nN
}
\]

of inclusion.

To see this, symmetry implies every population unit has the same
selection probability.

Since exactly \(n\) units are selected,

\[
N\pi=n,
\]

where \(\pi\) is the common inclusion probability.

Thus,

\[
\boxed{
\pi=\frac nN.
}
\]

%%%%%%%%%%%%%%%%%%%%%%%%%%%%%%%%%%%%%%%%%%%%%%%%%%%%%%%%%%%%
\subsection{Sampling With Replacement}
\label{subsec:sampling-with-replacement}
%%%%%%%%%%%%%%%%%%%%%%%%%%%%%%%%%%%%%%%%%%%%%%%%%%%%%%%%%%%%

Sampling with replacement means that after a unit is selected, it is
returned to the population before the next draw.

Thus the same unit may appear more than once.

If each draw selects uniformly from \(N\) units, then each draw has
probability

\[
\boxed{
\frac1N
}
\]

for every unit.

The draws are independent.

%%%%%%%%%%%%%%%%%%%%%%%%%%%%%%%%%%%%%%%%%%%%%%%%%%%%%%%%%%%%
\subsection{Sampling Without Replacement}
\label{subsec:sampling-without-replacement}
%%%%%%%%%%%%%%%%%%%%%%%%%%%%%%%%%%%%%%%%%%%%%%%%%%%%%%%%%%%%

Sampling without replacement means that once a unit is selected, it
cannot be selected again.

The draws are therefore dependent.

However, in an SRS, all unordered subsets of the specified size are
equally likely.

This is the standard interpretation of a simple random sample from a
finite population.

%%%%%%%%%%%%%%%%%%%%%%%%%%%%%%%%%%%%%%%%%%%%%%%%%%%%%%%%%%%%
\subsection{With Versus Without Replacement}
\label{subsec:with-versus-without-replacement}
%%%%%%%%%%%%%%%%%%%%%%%%%%%%%%%%%%%%%%%%%%%%%%%%%%%%%%%%%%%%

With replacement:

\[
\boxed{
\text{draws are independent}.
}
\]

Without replacement:

\[
\boxed{
\text{draws are dependent}.
}
\]

When the sample fraction

\[
\frac nN
\]

is very small, sampling without replacement may behave approximately like
independent sampling.

%%%%%%%%%%%%%%%%%%%%%%%%%%%%%%%%%%%%%%%%%%%%%%%%%%%%%%%%%%%%
\subsection{Finite-Population Correction}
\label{subsec:finite-population-correction}
%%%%%%%%%%%%%%%%%%%%%%%%%%%%%%%%%%%%%%%%%%%%%%%%%%%%%%%%%%%%

Sampling without replacement reduces variability because selected units
cannot reappear.

For an SRS from a finite population, the standard error of a sample mean
contains the finite-population correction

\[
\boxed{
\sqrt{
\frac{N-n}{N-1}
}.
}
\]

Thus a common variance formula is

\[
\boxed{
\operatorname{Var}(\overline{X})
=
\frac{\sigma^2}{n}
\frac{N-n}{N-1},
}
\]

where \(\sigma^2\) denotes the finite-population variance under the
corresponding denominator convention.

%%%%%%%%%%%%%%%%%%%%%%%%%%%%%%%%%%%%%%%%%%%%%%%%%%%%%%%%%%%%
\subsection{Sampling Fraction}
\label{subsec:sampling-fraction}
%%%%%%%%%%%%%%%%%%%%%%%%%%%%%%%%%%%%%%%%%%%%%%%%%%%%%%%%%%%%

The sampling fraction is

\[
\boxed{
f=\frac nN.
}
\]

If \(f\) is small, the finite-population correction is close to \(1\).

If \(f\) is large, sampling without replacement substantially reduces
sampling variance.

When

\[
n=N,
\]

the entire population is observed and sampling variance becomes zero.

%%%%%%%%%%%%%%%%%%%%%%%%%%%%%%%%%%%%%%%%%%%%%%%%%%%%%%%%%%%%
\subsection{Systematic Sampling}
\label{subsec:systematic-sampling}
%%%%%%%%%%%%%%%%%%%%%%%%%%%%%%%%%%%%%%%%%%%%%%%%%%%%%%%%%%%%

Systematic sampling selects units at regular intervals after a random
starting position.

For example, if approximately one in every \(k\) units is needed:

\begin{enumerate}
    \item randomly choose a starting position among the first \(k\)
          units;
    \item select every \(k\)-th unit afterward.
\end{enumerate}

Systematic sampling can be efficient when the population list is not
ordered in a problematic periodic pattern.

%%%%%%%%%%%%%%%%%%%%%%%%%%%%%%%%%%%%%%%%%%%%%%%%%%%%%%%%%%%%
\subsection{Potential Problem with Systematic Sampling}
\label{subsec:systematic-sampling-periodicity}
%%%%%%%%%%%%%%%%%%%%%%%%%%%%%%%%%%%%%%%%%%%%%%%%%%%%%%%%%%%%

If the population list contains periodic structure aligned with the
sampling interval, systematic sampling can produce bias or poor
representation.

For example, selecting every seventh record may be problematic if the
records repeat in a seven-unit cycle.

Thus the ordering of the sampling frame matters.

%%%%%%%%%%%%%%%%%%%%%%%%%%%%%%%%%%%%%%%%%%%%%%%%%%%%%%%%%%%%
\subsection{Stratified Sampling}
\label{subsec:stratified-sampling}
%%%%%%%%%%%%%%%%%%%%%%%%%%%%%%%%%%%%%%%%%%%%%%%%%%%%%%%%%%%%

\begin{definitionbox}{Stratified Sampling}
In stratified sampling, the population is divided into nonoverlapping
groups called strata, and samples are selected separately within the
strata.
\end{definitionbox}

If the strata are

\[
1,\ldots,H,
\]

then

\[
\boxed{
N
=
\sum_{h=1}^{H}N_h,
}
\]

where \(N_h\) is the population size of stratum \(h\).

%%%%%%%%%%%%%%%%%%%%%%%%%%%%%%%%%%%%%%%%%%%%%%%%%%%%%%%%%%%%
\subsection{Why Stratify?}
\label{subsec:why-stratify}
%%%%%%%%%%%%%%%%%%%%%%%%%%%%%%%%%%%%%%%%%%%%%%%%%%%%%%%%%%%%

Stratification can:

\[
\boxed{
\text{guarantee representation of important subgroups},
}
\]

\[
\boxed{
\text{increase precision},
}
\]

and

\[
\boxed{
\text{support subgroup estimates}.
}
\]

It is especially effective when units within each stratum are relatively
similar with respect to the variable of interest.

%%%%%%%%%%%%%%%%%%%%%%%%%%%%%%%%%%%%%%%%%%%%%%%%%%%%%%%%%%%%
\subsection{Proportional Allocation}
\label{subsec:proportional-allocation}
%%%%%%%%%%%%%%%%%%%%%%%%%%%%%%%%%%%%%%%%%%%%%%%%%%%%%%%%%%%%

Under proportional allocation,

\[
\boxed{
n_h
=
n
\frac{N_h}{N}.
}
\]

Thus each stratum receives approximately the same sampling fraction.

If the sample is self-weighting, every population unit may then have the
same overall inclusion probability.

%%%%%%%%%%%%%%%%%%%%%%%%%%%%%%%%%%%%%%%%%%%%%%%%%%%%%%%%%%%%
\subsection{Worked Example: Proportional Stratification}
\label{subsec:worked-proportional-stratification}
%%%%%%%%%%%%%%%%%%%%%%%%%%%%%%%%%%%%%%%%%%%%%%%%%%%%%%%%%%%%

\begin{workedexamplebox}{Sampling from Three Strata}

Suppose a population contains

\[
N=1000
\]

units divided into strata of sizes

\[
N_1=500,
\qquad
N_2=300,
\qquad
N_3=200.
\]

A proportional stratified sample of size

\[
n=100
\]

uses

\[
n_1=50,
\qquad
n_2=30,
\qquad
n_3=20.
\]

\begin{answerbox}
\[
\boxed{
(n_1,n_2,n_3)
=
(50,30,20).
}
\]
\end{answerbox}

\end{workedexamplebox}

%%%%%%%%%%%%%%%%%%%%%%%%%%%%%%%%%%%%%%%%%%%%%%%%%%%%%%%%%%%%
\subsection{Disproportionate Stratified Sampling}
\label{subsec:disproportionate-stratified-sampling}
%%%%%%%%%%%%%%%%%%%%%%%%%%%%%%%%%%%%%%%%%%%%%%%%%%%%%%%%%%%%

A study may intentionally oversample small or scientifically important
strata.

Then

\[
\frac{n_h}{n}
\]

need not equal

\[
\frac{N_h}{N}.
\]

Such a sample can improve subgroup precision but usually requires
weights when estimating population-wide quantities.

%%%%%%%%%%%%%%%%%%%%%%%%%%%%%%%%%%%%%%%%%%%%%%%%%%%%%%%%%%%%
\subsection{Neyman Allocation Preview}
\label{subsec:neyman-allocation-preview}
%%%%%%%%%%%%%%%%%%%%%%%%%%%%%%%%%%%%%%%%%%%%%%%%%%%%%%%%%%%%

If costs are similar across strata and within-stratum standard deviations
are known or estimated, Neyman allocation uses

\[
\boxed{
n_h
\propto
N_hS_h,
}
\]

where \(S_h\) measures within-stratum variability.

More variable strata receive larger samples.

This can reduce the variance of a stratified estimator.

%%%%%%%%%%%%%%%%%%%%%%%%%%%%%%%%%%%%%%%%%%%%%%%%%%%%%%%%%%%%
\subsection{Cluster Sampling}
\label{subsec:cluster-sampling}
%%%%%%%%%%%%%%%%%%%%%%%%%%%%%%%%%%%%%%%%%%%%%%%%%%%%%%%%%%%%

\begin{definitionbox}{Cluster Sampling}
Cluster sampling divides the population into naturally occurring groups
called clusters and randomly selects some clusters for observation.
\end{definitionbox}

Examples of clusters include:

\[
\text{schools},
\]

\[
\text{city blocks},
\]

\[
\text{households},
\]

or

\[
\text{hospitals}.
\]

%%%%%%%%%%%%%%%%%%%%%%%%%%%%%%%%%%%%%%%%%%%%%%%%%%%%%%%%%%%%
\subsection{One-Stage Cluster Sampling}
\label{subsec:one-stage-cluster-sampling}
%%%%%%%%%%%%%%%%%%%%%%%%%%%%%%%%%%%%%%%%%%%%%%%%%%%%%%%%%%%%

In one-stage cluster sampling:

\begin{enumerate}
    \item randomly select clusters;
    \item observe every eligible unit inside selected clusters.
\end{enumerate}

This can greatly reduce travel and administrative costs.

%%%%%%%%%%%%%%%%%%%%%%%%%%%%%%%%%%%%%%%%%%%%%%%%%%%%%%%%%%%%
\subsection{Two-Stage Cluster Sampling}
\label{subsec:two-stage-cluster-sampling}
%%%%%%%%%%%%%%%%%%%%%%%%%%%%%%%%%%%%%%%%%%%%%%%%%%%%%%%%%%%%

In two-stage sampling:

\begin{enumerate}
    \item randomly select clusters;
    \item randomly sample units within each selected cluster.
\end{enumerate}

This is a common form of multistage sampling.

Large national surveys often use several stages.

%%%%%%%%%%%%%%%%%%%%%%%%%%%%%%%%%%%%%%%%%%%%%%%%%%%%%%%%%%%%
\subsection{Strata Versus Clusters}
\label{subsec:strata-versus-clusters}
%%%%%%%%%%%%%%%%%%%%%%%%%%%%%%%%%%%%%%%%%%%%%%%%%%%%%%%%%%%%

Strata and clusters serve different purposes.

Good strata are often internally homogeneous but different from one
another.

Good clusters for economical sampling may contain diverse units so that
each cluster resembles a small version of the full population.

Thus:

\[
\boxed{
\text{stratification}
\rightarrow
\text{sample from every stratum},
}
\]

while

\[
\boxed{
\text{cluster sampling}
\rightarrow
\text{sample only selected clusters}.
}
\]

%%%%%%%%%%%%%%%%%%%%%%%%%%%%%%%%%%%%%%%%%%%%%%%%%%%%%%%%%%%%
\subsection{Cluster Correlation}
\label{subsec:cluster-correlation}
%%%%%%%%%%%%%%%%%%%%%%%%%%%%%%%%%%%%%%%%%%%%%%%%%%%%%%%%%%%%

Units inside the same cluster may be more similar than units from
different clusters.

This is called intracluster correlation.

If cluster members are highly similar, a sample of many units from only
a few clusters may contain less information than an SRS of the same
size.

Thus cluster sampling can increase variance even while reducing cost.

%%%%%%%%%%%%%%%%%%%%%%%%%%%%%%%%%%%%%%%%%%%%%%%%%%%%%%%%%%%%
\subsection{Design Effect Preview}
\label{subsec:design-effect-preview}
%%%%%%%%%%%%%%%%%%%%%%%%%%%%%%%%%%%%%%%%%%%%%%%%%%%%%%%%%%%%

The design effect compares variance under a complex sample design with
variance under an SRS of the same nominal sample size.

One common definition is

\[
\boxed{
DEFF
=
\frac{
\operatorname{Var}_{\text{design}}(\widehat{\theta})
}{
\operatorname{Var}_{\text{SRS}}(\widehat{\theta})
}.
}
\]

If

\[
DEFF>1,
\]

the design has greater variance than the comparable SRS.

%%%%%%%%%%%%%%%%%%%%%%%%%%%%%%%%%%%%%%%%%%%%%%%%%%%%%%%%%%%%
\subsection{Multistage Sampling}
\label{subsec:multistage-sampling}
%%%%%%%%%%%%%%%%%%%%%%%%%%%%%%%%%%%%%%%%%%%%%%%%%%%%%%%%%%%%

Multistage sampling selects units in several randomized stages.

For example:

\[
\boxed{
\text{counties}
\to
\text{schools}
\to
\text{classrooms}
\to
\text{students}.
}
\]

Different sampling probabilities may occur at each stage.

The final inclusion probability is determined by the product of relevant
conditional selection probabilities.

%%%%%%%%%%%%%%%%%%%%%%%%%%%%%%%%%%%%%%%%%%%%%%%%%%%%%%%%%%%%
\subsection{Unequal-Probability Sampling}
\label{subsec:unequal-probability-sampling}
%%%%%%%%%%%%%%%%%%%%%%%%%%%%%%%%%%%%%%%%%%%%%%%%%%%%%%%%%%%%

Not every probability design gives every unit the same inclusion
probability.

Let

\[
\pi_i
=
P(\text{unit }i\text{ is included}).
\]

If the \(\pi_i\) differ, then unweighted sample averages may not estimate
population averages correctly.

Sampling weights can adjust for unequal selection probabilities.

%%%%%%%%%%%%%%%%%%%%%%%%%%%%%%%%%%%%%%%%%%%%%%%%%%%%%%%%%%%%
\subsection{Sampling Weights}
\label{subsec:sampling-weights}
%%%%%%%%%%%%%%%%%%%%%%%%%%%%%%%%%%%%%%%%%%%%%%%%%%%%%%%%%%%%

A basic design weight is

\[
\boxed{
w_i
=
\frac1{\pi_i}.
}
\]

A unit with a small inclusion probability receives a larger weight.

The intuition is that such a sampled unit represents more population
units.

%%%%%%%%%%%%%%%%%%%%%%%%%%%%%%%%%%%%%%%%%%%%%%%%%%%%%%%%%%%%
\subsection{Horvitz--Thompson Estimator Preview}
\label{subsec:horvitz-thompson-preview}
%%%%%%%%%%%%%%%%%%%%%%%%%%%%%%%%%%%%%%%%%%%%%%%%%%%%%%%%%%%%

Suppose population values are \(y_i\), and \(I_i\) indicates whether unit
\(i\) is included.

A classic unbiased estimator of the population total is

\[
\boxed{
\widehat{Y}_{HT}
=
\sum_{i=1}^{N}
\frac{I_iy_i}{\pi_i}.
}
\]

Since

\[
\mathbb{E}[I_i]=\pi_i,
\]

we obtain

\[
\mathbb{E}
\left[
\frac{I_iy_i}{\pi_i}
\right]
=
y_i.
\]

Therefore,

\[
\boxed{
\mathbb{E}[\widehat{Y}_{HT}]
=
\sum_{i=1}^{N}y_i.
}
\]

%%%%%%%%%%%%%%%%%%%%%%%%%%%%%%%%%%%%%%%%%%%%%%%%%%%%%%%%%%%%
\subsection{Sampling Error}
\label{subsec:sampling-error}
%%%%%%%%%%%%%%%%%%%%%%%%%%%%%%%%%%%%%%%%%%%%%%%%%%%%%%%%%%%%

Sampling error is the difference between a sample statistic and its
target population parameter caused by observing only a sample.

For example,

\[
\boxed{
\overline{X}-\mu
}
\]

is sampling error for a sample mean.

Even a perfectly designed probability sample has random sampling error.

Sampling distributions quantify this variability.

%%%%%%%%%%%%%%%%%%%%%%%%%%%%%%%%%%%%%%%%%%%%%%%%%%%%%%%%%%%%
\subsection{Nonsampling Error}
\label{subsec:nonsampling-error}
%%%%%%%%%%%%%%%%%%%%%%%%%%%%%%%%%%%%%%%%%%%%%%%%%%%%%%%%%%%%

Nonsampling error includes errors not caused merely by selecting a
sample.

Examples include:

\[
\boxed{
\text{undercoverage},
}
\]

\[
\boxed{
\text{nonresponse},
}
\]

\[
\boxed{
\text{measurement error},
}
\]

\[
\boxed{
\text{data-entry errors},
}
\]

and

\[
\boxed{
\text{questionnaire effects}.
}
\]

Increasing sample size does not automatically eliminate these errors.

%%%%%%%%%%%%%%%%%%%%%%%%%%%%%%%%%%%%%%%%%%%%%%%%%%%%%%%%%%%%
\subsection{Bias}
\label{subsec:bias-data-collection}
%%%%%%%%%%%%%%%%%%%%%%%%%%%%%%%%%%%%%%%%%%%%%%%%%%%%%%%%%%%%

Bias is a systematic tendency for a procedure to favor certain outcomes
or estimates.

A biased design may produce estimates that repeatedly miss the target in
the same direction.

Bias differs from random sampling variability.

A large sample can have:

\[
\boxed{
\text{small random error}
}
\]

but still

\[
\boxed{
\text{large systematic bias}.
}
\]

%%%%%%%%%%%%%%%%%%%%%%%%%%%%%%%%%%%%%%%%%%%%%%%%%%%%%%%%%%%%
\subsection{Undercoverage}
\label{subsec:undercoverage}
%%%%%%%%%%%%%%%%%%%%%%%%%%%%%%%%%%%%%%%%%%%%%%%%%%%%%%%%%%%%

Undercoverage occurs when some members of the target population are
missing or inadequately represented in the sampling frame.

For example, a survey based only on landline telephone numbers may omit
people without landlines.

If omitted groups differ systematically from included groups, estimates
may be biased.

%%%%%%%%%%%%%%%%%%%%%%%%%%%%%%%%%%%%%%%%%%%%%%%%%%%%%%%%%%%%
\subsection{Selection Bias}
\label{subsec:selection-bias}
%%%%%%%%%%%%%%%%%%%%%%%%%%%%%%%%%%%%%%%%%%%%%%%%%%%%%%%%%%%%

Selection bias occurs when the mechanism determining inclusion in the
observed sample is related to the variables being studied.

For example, a health study based only on volunteers may attract people
who are unusually health-conscious.

The sample can then differ systematically from the target population.

%%%%%%%%%%%%%%%%%%%%%%%%%%%%%%%%%%%%%%%%%%%%%%%%%%%%%%%%%%%%
\subsection{Nonresponse Bias}
\label{subsec:nonresponse-bias}
%%%%%%%%%%%%%%%%%%%%%%%%%%%%%%%%%%%%%%%%%%%%%%%%%%%%%%%%%%%%

Nonresponse occurs when selected units do not provide usable data.

Nonresponse bias arises when respondents differ systematically from
nonrespondents on variables related to the study.

A low response rate does not mathematically guarantee bias, but it can
increase concern about representativeness.

%%%%%%%%%%%%%%%%%%%%%%%%%%%%%%%%%%%%%%%%%%%%%%%%%%%%%%%%%%%%
\subsection{Response Bias}
\label{subsec:response-bias}
%%%%%%%%%%%%%%%%%%%%%%%%%%%%%%%%%%%%%%%%%%%%%%%%%%%%%%%%%%%%

Response bias occurs when recorded answers systematically differ from
the underlying truth.

Possible causes include:

\[
\boxed{
\text{social desirability},
}
\]

\[
\boxed{
\text{memory limitations},
}
\]

\[
\boxed{
\text{interviewer effects},
}
\]

and

\[
\boxed{
\text{question wording}.
}
\]

%%%%%%%%%%%%%%%%%%%%%%%%%%%%%%%%%%%%%%%%%%%%%%%%%%%%%%%%%%%%
\subsection{Leading Questions}
\label{subsec:leading-questions}
%%%%%%%%%%%%%%%%%%%%%%%%%%%%%%%%%%%%%%%%%%%%%%%%%%%%%%%%%%%%

A leading question encourages a particular answer.

For example, wording such as

\[
\text{``Do you agree that the wasteful program should be reduced?''}
\]

contains emotionally loaded language.

A more neutral form would describe the policy without presupposing that
it is wasteful.

Survey wording is part of measurement design.

%%%%%%%%%%%%%%%%%%%%%%%%%%%%%%%%%%%%%%%%%%%%%%%%%%%%%%%%%%%%
\subsection{Question Order Effects}
\label{subsec:question-order-effects}
%%%%%%%%%%%%%%%%%%%%%%%%%%%%%%%%%%%%%%%%%%%%%%%%%%%%%%%%%%%%

Responses may depend on questions asked earlier in the survey.

An earlier question can influence:

\[
\text{memory},
\]

\[
\text{interpretation},
\]

or

\[
\text{emotional context}.
\]

Randomizing question order may sometimes help diagnose or reduce such
effects.

%%%%%%%%%%%%%%%%%%%%%%%%%%%%%%%%%%%%%%%%%%%%%%%%%%%%%%%%%%%%
\subsection{Measurement Error}
\label{subsec:measurement-error}
%%%%%%%%%%%%%%%%%%%%%%%%%%%%%%%%%%%%%%%%%%%%%%%%%%%%%%%%%%%%

Suppose the true quantity is

\[
X
\]

but the observed quantity is

\[
Y.
\]

A simple measurement-error model is

\[
\boxed{
Y=X+\varepsilon,
}
\]

where

\[
\varepsilon
\]

represents measurement noise.

Measurement error may arise from instruments, observers, coding, or
reporting behavior.

%%%%%%%%%%%%%%%%%%%%%%%%%%%%%%%%%%%%%%%%%%%%%%%%%%%%%%%%%%%%
\subsection{Validity and Reliability}
\label{subsec:validity-reliability}
%%%%%%%%%%%%%%%%%%%%%%%%%%%%%%%%%%%%%%%%%%%%%%%%%%%%%%%%%%%%

A measurement is valid if it measures the intended concept.

A measurement is reliable if it produces consistent results under
repeated comparable conditions.

A measurement can be reliable but invalid.

For example, a miscalibrated scale may consistently report values that
are too large.

%%%%%%%%%%%%%%%%%%%%%%%%%%%%%%%%%%%%%%%%%%%%%%%%%%%%%%%%%%%%
\subsection{Observational Studies}
\label{subsec:observational-studies}
%%%%%%%%%%%%%%%%%%%%%%%%%%%%%%%%%%%%%%%%%%%%%%%%%%%%%%%%%%%%

\begin{definitionbox}{Observational Study}
An observational study records variables without deliberately assigning
treatments to experimental units.
\end{definitionbox}

Researchers observe naturally occurring differences.

Examples include:

\[
\text{survey studies},
\]

\[
\text{cohort studies},
\]

and

\[
\text{many epidemiological studies}.
\]

%%%%%%%%%%%%%%%%%%%%%%%%%%%%%%%%%%%%%%%%%%%%%%%%%%%%%%%%%%%%
\subsection{Experiments}
\label{subsec:experiments}
%%%%%%%%%%%%%%%%%%%%%%%%%%%%%%%%%%%%%%%%%%%%%%%%%%%%%%%%%%%%

\begin{definitionbox}{Experiment}
An experiment deliberately imposes one or more treatments on
experimental units and measures resulting responses.
\end{definitionbox}

Well-designed randomized experiments provide stronger evidence for
causal effects than ordinary observational studies.

%%%%%%%%%%%%%%%%%%%%%%%%%%%%%%%%%%%%%%%%%%%%%%%%%%%%%%%%%%%%
\subsection{Experimental Units}
\label{subsec:experimental-units}
%%%%%%%%%%%%%%%%%%%%%%%%%%%%%%%%%%%%%%%%%%%%%%%%%%%%%%%%%%%%

Experimental units are the objects to which treatments are assigned.

Depending on the study, they may be:

\[
\text{people},
\]

\[
\text{plots of land},
\]

\[
\text{machines},
\]

\[
\text{laboratory samples},
\]

or

\[
\text{schools}.
\]

When the units are human, they are often called subjects or participants.

%%%%%%%%%%%%%%%%%%%%%%%%%%%%%%%%%%%%%%%%%%%%%%%%%%%%%%%%%%%%
\subsection{Treatment}
\label{subsec:treatment}
%%%%%%%%%%%%%%%%%%%%%%%%%%%%%%%%%%%%%%%%%%%%%%%%%%%%%%%%%%%%

A treatment is a specific experimental condition imposed on an
experimental unit.

If an experiment varies several factors, a treatment may be a particular
combination of factor levels.

For example:

\[
\text{fertilizer type}
\]

and

\[
\text{water level}
\]

may jointly determine a treatment.

%%%%%%%%%%%%%%%%%%%%%%%%%%%%%%%%%%%%%%%%%%%%%%%%%%%%%%%%%%%%
\subsection{Response Variable}
\label{subsec:response-variable}
%%%%%%%%%%%%%%%%%%%%%%%%%%%%%%%%%%%%%%%%%%%%%%%%%%%%%%%%%%%%

The response variable is the outcome measured to evaluate the effect of
a treatment or explanatory variable.

Examples include:

\[
\text{blood pressure},
\]

\[
\text{test score},
\]

\[
\text{crop yield},
\]

or

\[
\text{machine lifetime}.
\]

%%%%%%%%%%%%%%%%%%%%%%%%%%%%%%%%%%%%%%%%%%%%%%%%%%%%%%%%%%%%
\subsection{Explanatory Variable}
\label{subsec:explanatory-variable}
%%%%%%%%%%%%%%%%%%%%%%%%%%%%%%%%%%%%%%%%%%%%%%%%%%%%%%%%%%%%

An explanatory variable is a variable used to explain or predict changes
in the response.

In an experiment, explanatory variables may be controlled factors.

In an observational study, they may be naturally occurring
characteristics.

%%%%%%%%%%%%%%%%%%%%%%%%%%%%%%%%%%%%%%%%%%%%%%%%%%%%%%%%%%%%
\subsection{Confounding}
\label{subsec:confounding}
%%%%%%%%%%%%%%%%%%%%%%%%%%%%%%%%%%%%%%%%%%%%%%%%%%%%%%%%%%%%

Confounding occurs when the effects of two or more explanatory variables
cannot be separated clearly.

Suppose an observational study finds an association between

\[
X
\]

and

\[
Y.
\]

A third variable

\[
Z
\]

may influence both.

Then the observed \(X\)-\(Y\) relationship may partly or entirely reflect
the effect of \(Z\).

%%%%%%%%%%%%%%%%%%%%%%%%%%%%%%%%%%%%%%%%%%%%%%%%%%%%%%%%%%%%
\subsection{Worked Example: Confounding}
\label{subsec:worked-confounding}
%%%%%%%%%%%%%%%%%%%%%%%%%%%%%%%%%%%%%%%%%%%%%%%%%%%%%%%%%%%%

\begin{workedexamplebox}{Exercise and Health}

Suppose an observational study finds that people who exercise more have
better cardiovascular outcomes.

Age may be associated with both exercise behavior and cardiovascular
health.

Therefore age may act as a confounding variable.

\begin{answerbox}
\[
\boxed{
\text{Association between exercise and health does not by itself prove
causation.}
}
\]
\end{answerbox}

Additional design or adjustment is needed.

\end{workedexamplebox}

%%%%%%%%%%%%%%%%%%%%%%%%%%%%%%%%%%%%%%%%%%%%%%%%%%%%%%%%%%%%
\subsection{Random Assignment}
\label{subsec:random-assignment}
%%%%%%%%%%%%%%%%%%%%%%%%%%%%%%%%%%%%%%%%%%%%%%%%%%%%%%%%%%%%

Random assignment uses chance to determine which experimental units
receive which treatments.

The goal is to make treatment groups comparable with respect to both
measured and unmeasured characteristics.

Random assignment supports causal inference by breaking systematic links
between treatment and potential confounders.

%%%%%%%%%%%%%%%%%%%%%%%%%%%%%%%%%%%%%%%%%%%%%%%%%%%%%%%%%%%%
\subsection{Random Sampling Versus Random Assignment}
\label{subsec:random-sampling-versus-assignment}
%%%%%%%%%%%%%%%%%%%%%%%%%%%%%%%%%%%%%%%%%%%%%%%%%%%%%%%%%%%%

These concepts have different purposes.

Random sampling addresses:

\[
\boxed{
\text{Can results be generalized to a population?}
}
\]

Random assignment addresses:

\[
\boxed{
\text{Can treatment differences be interpreted causally?}
}
\]

A study may have one without the other.

%%%%%%%%%%%%%%%%%%%%%%%%%%%%%%%%%%%%%%%%%%%%%%%%%%%%%%%%%%%%
\subsection{Four Design Combinations}
\label{subsec:four-randomization-combinations}
%%%%%%%%%%%%%%%%%%%%%%%%%%%%%%%%%%%%%%%%%%%%%%%%%%%%%%%%%%%%

A useful conceptual table is

\[
\begin{array}{c|c|c}
&
\text{Random Assignment}
&
\text{No Random Assignment}
\\
\hline
\text{Random Sampling}
&
\begin{array}{c}
\text{population inference}\\
+\text{ stronger causal inference}
\end{array}
&
\begin{array}{c}
\text{population association}\\
\text{but limited causal claim}
\end{array}
\\[4mm]
\text{No Random Sampling}
&
\begin{array}{c}
\text{causal inference for}\\
\text{studied units/population}\\
\text{under assumptions}
\end{array}
&
\begin{array}{c}
\text{limited generalization}\\
\text{and limited causality}
\end{array}
\end{array}
\]

The exact strength of conclusions also depends on implementation and
study context.

%%%%%%%%%%%%%%%%%%%%%%%%%%%%%%%%%%%%%%%%%%%%%%%%%%%%%%%%%%%%
\subsection{Control Group}
\label{subsec:control-group}
%%%%%%%%%%%%%%%%%%%%%%%%%%%%%%%%%%%%%%%%%%%%%%%%%%%%%%%%%%%%

A control group provides a reference condition against which one or more
treatments are compared.

The control may receive:

\[
\text{no intervention},
\]

\[
\text{standard treatment},
\]

or

\[
\text{a placebo}.
\]

The appropriate control depends on the scientific question.

%%%%%%%%%%%%%%%%%%%%%%%%%%%%%%%%%%%%%%%%%%%%%%%%%%%%%%%%%%%%
\subsection{Placebo}
\label{subsec:placebo}
%%%%%%%%%%%%%%%%%%%%%%%%%%%%%%%%%%%%%%%%%%%%%%%%%%%%%%%%%%%%

A placebo is an inactive or comparison treatment designed to resemble
the active treatment.

Placebos can help separate the treatment's specific effect from changes
associated with expectations, attention, or study participation.

Not every experiment can or should use a placebo.

%%%%%%%%%%%%%%%%%%%%%%%%%%%%%%%%%%%%%%%%%%%%%%%%%%%%%%%%%%%%
\subsection{Blinding}
\label{subsec:blinding}
%%%%%%%%%%%%%%%%%%%%%%%%%%%%%%%%%%%%%%%%%%%%%%%%%%%%%%%%%%%%

Blinding means withholding treatment information from some participants
in the study.

A participant may be blinded to treatment assignment.

An evaluator may also be blinded.

Blinding can reduce behavioral and measurement biases.

%%%%%%%%%%%%%%%%%%%%%%%%%%%%%%%%%%%%%%%%%%%%%%%%%%%%%%%%%%%%
\subsection{Double Blinding}
\label{subsec:double-blinding}
%%%%%%%%%%%%%%%%%%%%%%%%%%%%%%%%%%%%%%%%%%%%%%%%%%%%%%%%%%%%

In a double-blind study, both participants and key evaluators or
treatment administrators are unaware of treatment assignments, when
feasible.

This can reduce:

\[
\boxed{
\text{participant expectation effects}
}
\]

and

\[
\boxed{
\text{observer expectation effects}.
}
\]

%%%%%%%%%%%%%%%%%%%%%%%%%%%%%%%%%%%%%%%%%%%%%%%%%%%%%%%%%%%%
\subsection{Replication}
\label{subsec:replication}
%%%%%%%%%%%%%%%%%%%%%%%%%%%%%%%%%%%%%%%%%%%%%%%%%%%%%%%%%%%%

Replication means applying treatments to multiple experimental units.

Replication provides information about natural variability.

A treatment applied to only one unit cannot be separated from the unique
characteristics of that unit.

%%%%%%%%%%%%%%%%%%%%%%%%%%%%%%%%%%%%%%%%%%%%%%%%%%%%%%%%%%%%
\subsection{Randomization}
\label{subsec:randomization-experimental-design}
%%%%%%%%%%%%%%%%%%%%%%%%%%%%%%%%%%%%%%%%%%%%%%%%%%%%%%%%%%%%

Randomization protects against systematic allocation bias.

It also provides a probability basis for many inferential procedures.

Random assignment does not guarantee perfectly balanced groups in every
finite experiment.

Instead, it makes systematic imbalance less likely and gives a known
random mechanism for treatment allocation.

%%%%%%%%%%%%%%%%%%%%%%%%%%%%%%%%%%%%%%%%%%%%%%%%%%%%%%%%%%%%
\subsection{Blocking}
\label{subsec:blocking}
%%%%%%%%%%%%%%%%%%%%%%%%%%%%%%%%%%%%%%%%%%%%%%%%%%%%%%%%%%%%

Blocking groups similar experimental units before randomization.

Treatments are then randomized within each block.

If an important source of variation is known in advance, blocking can
reduce unexplained variability.

%%%%%%%%%%%%%%%%%%%%%%%%%%%%%%%%%%%%%%%%%%%%%%%%%%%%%%%%%%%%
\subsection{Worked Example: Blocking}
\label{subsec:worked-blocking}
%%%%%%%%%%%%%%%%%%%%%%%%%%%%%%%%%%%%%%%%%%%%%%%%%%%%%%%%%%%%

\begin{workedexamplebox}{Blocking by Initial Skill Level}

Suppose two teaching methods are compared.

Students' prior mathematics preparation strongly affects final scores.

Students may first be grouped into blocks such as:

\[
\text{low preparation},
\]

\[
\text{medium preparation},
\]

and

\[
\text{high preparation}.
\]

Within each block, students are randomly assigned to the two methods.

\begin{answerbox}
\[
\boxed{
\text{Blocking helps separate treatment effects from preparation-level
variation.}
}
\]
\end{answerbox}

\end{workedexamplebox}

%%%%%%%%%%%%%%%%%%%%%%%%%%%%%%%%%%%%%%%%%%%%%%%%%%%%%%%%%%%%
\subsection{Matched-Pairs Design}
\label{subsec:matched-pairs-design}
%%%%%%%%%%%%%%%%%%%%%%%%%%%%%%%%%%%%%%%%%%%%%%%%%%%%%%%%%%%%

A matched-pairs design is a special blocking design.

Two common forms are:

\begin{enumerate}
    \item pair similar experimental units and assign different treatments
          within each pair;
    \item measure the same unit under two conditions.
\end{enumerate}

The analysis then focuses on within-pair differences.

%%%%%%%%%%%%%%%%%%%%%%%%%%%%%%%%%%%%%%%%%%%%%%%%%%%%%%%%%%%%
\subsection{Before-and-After Studies}
\label{subsec:before-after-studies}
%%%%%%%%%%%%%%%%%%%%%%%%%%%%%%%%%%%%%%%%%%%%%%%%%%%%%%%%%%%%

A before-and-after design measures each subject before and after an
intervention.

Define

\[
\boxed{
D_i
=
Y_{i,\text{after}}
-
Y_{i,\text{before}}.
}
\]

The resulting data consist of paired differences

\[
D_1,\ldots,D_n.
\]

This can reduce variation caused by stable subject-to-subject
differences.

%%%%%%%%%%%%%%%%%%%%%%%%%%%%%%%%%%%%%%%%%%%%%%%%%%%%%%%%%%%%
\subsection{Factor}
\label{subsec:factor-experiment}
%%%%%%%%%%%%%%%%%%%%%%%%%%%%%%%%%%%%%%%%%%%%%%%%%%%%%%%%%%%%

A factor is an explanatory variable deliberately varied in an
experiment.

For example:

\[
\text{temperature}
\]

may be a factor with levels

\[
10^\circ,
\quad
20^\circ,
\quad
30^\circ.
\]

Factorial experimental designs are developed later in
Section~\ref{sec:experimental-design}.

%%%%%%%%%%%%%%%%%%%%%%%%%%%%%%%%%%%%%%%%%%%%%%%%%%%%%%%%%%%%
\subsection{Cross-Sectional Studies}
\label{subsec:cross-sectional-studies}
%%%%%%%%%%%%%%%%%%%%%%%%%%%%%%%%%%%%%%%%%%%%%%%%%%%%%%%%%%%%

A cross-sectional study measures units at approximately one point in
time.

It provides a snapshot of a population or system.

Examples include:

\[
\text{one-time surveys},
\]

or

\[
\text{measurements taken during one examination period}.
\]

%%%%%%%%%%%%%%%%%%%%%%%%%%%%%%%%%%%%%%%%%%%%%%%%%%%%%%%%%%%%
\subsection{Longitudinal Studies}
\label{subsec:longitudinal-studies}
%%%%%%%%%%%%%%%%%%%%%%%%%%%%%%%%%%%%%%%%%%%%%%%%%%%%%%%%%%%%

A longitudinal study follows units over time.

Repeated observations may be collected from the same participants.

This allows study of:

\[
\boxed{
\text{change},
}
\]

\[
\boxed{
\text{trajectories},
}
\]

and

\[
\boxed{
\text{within-subject variation}.
}
\]

Repeated observations on the same unit are generally dependent.

%%%%%%%%%%%%%%%%%%%%%%%%%%%%%%%%%%%%%%%%%%%%%%%%%%%%%%%%%%%%
\subsection{Prospective Studies}
\label{subsec:prospective-studies}
%%%%%%%%%%%%%%%%%%%%%%%%%%%%%%%%%%%%%%%%%%%%%%%%%%%%%%%%%%%%

A prospective study identifies participants and follows them forward in
time.

Exposures or characteristics may be recorded before future outcomes are
observed.

Prospective cohort studies are common in epidemiology.

%%%%%%%%%%%%%%%%%%%%%%%%%%%%%%%%%%%%%%%%%%%%%%%%%%%%%%%%%%%%
\subsection{Retrospective Studies}
\label{subsec:retrospective-studies}
%%%%%%%%%%%%%%%%%%%%%%%%%%%%%%%%%%%%%%%%%%%%%%%%%%%%%%%%%%%%

A retrospective study uses historical information about events that have
already occurred.

For example, investigators may identify individuals with and without a
condition and examine previous exposures.

Retrospective designs can be efficient but may suffer from incomplete
records or recall bias.

%%%%%%%%%%%%%%%%%%%%%%%%%%%%%%%%%%%%%%%%%%%%%%%%%%%%%%%%%%%%
\subsection{Case-Control Studies}
\label{subsec:case-control-studies}
%%%%%%%%%%%%%%%%%%%%%%%%%%%%%%%%%%%%%%%%%%%%%%%%%%%%%%%%%%%%

A case-control study compares:

\[
\boxed{
\text{cases}
=
\text{units with the outcome},
}
\]

and

\[
\boxed{
\text{controls}
=
\text{units without the outcome}.
}
\]

Investigators then compare previous exposures.

Case-control sampling does not directly estimate disease prevalence from
the sample fraction because sampling is conditioned on outcome status.

%%%%%%%%%%%%%%%%%%%%%%%%%%%%%%%%%%%%%%%%%%%%%%%%%%%%%%%%%%%%
\subsection{Cohort Studies}
\label{subsec:cohort-studies}
%%%%%%%%%%%%%%%%%%%%%%%%%%%%%%%%%%%%%%%%%%%%%%%%%%%%%%%%%%%%

A cohort study follows groups defined by exposure or other baseline
characteristics and compares subsequent outcomes.

If outcome risks can be directly observed, cohort studies can estimate
quantities such as:

\[
\boxed{
\text{risk},
}
\]

\[
\boxed{
\text{risk difference},
}
\]

and

\[
\boxed{
\text{relative risk}.
}
\]

%%%%%%%%%%%%%%%%%%%%%%%%%%%%%%%%%%%%%%%%%%%%%%%%%%%%%%%%%%%%
\subsection{Administrative Data}
\label{subsec:administrative-data}
%%%%%%%%%%%%%%%%%%%%%%%%%%%%%%%%%%%%%%%%%%%%%%%%%%%%%%%%%%%%

Administrative data are collected primarily for operational purposes
rather than research.

Examples include:

\[
\text{billing records},
\]

\[
\text{school enrollment records},
\]

\[
\text{tax records},
\]

and

\[
\text{service logs}.
\]

These datasets may be large but can contain variables defined for
administrative rather than scientific purposes.

%%%%%%%%%%%%%%%%%%%%%%%%%%%%%%%%%%%%%%%%%%%%%%%%%%%%%%%%%%%%
\subsection{Secondary Data}
\label{subsec:secondary-data}
%%%%%%%%%%%%%%%%%%%%%%%%%%%%%%%%%%%%%%%%%%%%%%%%%%%%%%%%%%%%

Secondary data are data collected previously by another person,
organization, or study.

Before using secondary data, investigators should understand:

\[
\text{original purpose},
\]

\[
\text{sampling design},
\]

\[
\text{variable definitions},
\]

\[
\text{missing-data patterns},
\]

and

\[
\text{possible changes across time}.
\]

%%%%%%%%%%%%%%%%%%%%%%%%%%%%%%%%%%%%%%%%%%%%%%%%%%%%%%%%%%%%
\subsection{Big Data Does Not Eliminate Bias}
\label{subsec:big-data-bias}
%%%%%%%%%%%%%%%%%%%%%%%%%%%%%%%%%%%%%%%%%%%%%%%%%%%%%%%%%%%%

A dataset with millions of observations may still be systematically
unrepresentative.

Large \(n\) reduces certain forms of random sampling error.

It does not automatically eliminate:

\[
\boxed{
\text{coverage bias},
}
\]

\[
\boxed{
\text{selection bias},
}
\]

\[
\boxed{
\text{measurement bias},
}
\]

or

\[
\boxed{
\text{confounding}.
}
\]

Thus sample quality is not determined only by sample size.

%%%%%%%%%%%%%%%%%%%%%%%%%%%%%%%%%%%%%%%%%%%%%%%%%%%%%%%%%%%%
\subsection{Design-Based Inference}
\label{subsec:design-based-inference}
%%%%%%%%%%%%%%%%%%%%%%%%%%%%%%%%%%%%%%%%%%%%%%%%%%%%%%%%%%%%

In design-based inference, the finite population values are treated as
fixed.

Randomness comes from the sampling design.

For example, a sample mean is random because different random samples
could have been selected.

This perspective is central in survey sampling.

%%%%%%%%%%%%%%%%%%%%%%%%%%%%%%%%%%%%%%%%%%%%%%%%%%%%%%%%%%%%
\subsection{Model-Based Inference}
\label{subsec:model-based-inference}
%%%%%%%%%%%%%%%%%%%%%%%%%%%%%%%%%%%%%%%%%%%%%%%%%%%%%%%%%%%%

In model-based inference, observations are treated as realizations from a
probability model.

For example,

\[
X_1,\ldots,X_n
\overset{\text{i.i.d.}}{\sim}
F_\theta.
\]

Randomness comes from the model generating the observations.

Much of classical mathematical statistics uses this perspective.

%%%%%%%%%%%%%%%%%%%%%%%%%%%%%%%%%%%%%%%%%%%%%%%%%%%%%%%%%%%%
\subsection{Design-Based Versus Model-Based Perspectives}
\label{subsec:design-model-comparison}
%%%%%%%%%%%%%%%%%%%%%%%%%%%%%%%%%%%%%%%%%%%%%%%%%%%%%%%%%%%%

The two approaches emphasize different sources of randomness.

Design-based:

\[
\boxed{
\text{random sample selection}.
}
\]

Model-based:

\[
\boxed{
\text{random data-generating mechanism}.
}
\]

Modern statistical practice may combine both perspectives.

%%%%%%%%%%%%%%%%%%%%%%%%%%%%%%%%%%%%%%%%%%%%%%%%%%%%%%%%%%%%
\subsection{Sampling Distribution Preview}
\label{subsec:sampling-distribution-preview}
%%%%%%%%%%%%%%%%%%%%%%%%%%%%%%%%%%%%%%%%%%%%%%%%%%%%%%%%%%%%

Suppose an SRS is repeatedly drawn and a statistic

\[
T
\]

is calculated each time.

The probability distribution of

\[
T
\]

over all possible samples is its sampling distribution.

For example,

\[
\boxed{
\overline{X}
}
\]

has a sampling distribution.

The next section develops sampling distributions systematically.

%%%%%%%%%%%%%%%%%%%%%%%%%%%%%%%%%%%%%%%%%%%%%%%%%%%%%%%%%%%%
\subsection{Standard Error Preview}
\label{subsec:standard-error-sampling-preview}
%%%%%%%%%%%%%%%%%%%%%%%%%%%%%%%%%%%%%%%%%%%%%%%%%%%%%%%%%%%%

The standard deviation of a statistic's sampling distribution is called
its standard error.

For an i.i.d. sample mean,

\[
\boxed{
\operatorname{SE}(\overline{X})
=
\frac{\sigma}{\sqrt{n}}.
}
\]

For finite-population sampling without replacement, the finite-population
correction may also appear.

%%%%%%%%%%%%%%%%%%%%%%%%%%%%%%%%%%%%%%%%%%%%%%%%%%%%%%%%%%%%
\subsection{Representative Samples}
\label{subsec:representative-samples}
%%%%%%%%%%%%%%%%%%%%%%%%%%%%%%%%%%%%%%%%%%%%%%%%%%%%%%%%%%%%

The phrase ``representative sample'' is often used informally.

A probability sample does not guarantee that every sample exactly
matches the population.

Rather, a known random selection design allows sampling variability to
be quantified and makes systematic selection bias less likely.

Representativeness is therefore a property to evaluate, not something
guaranteed merely by sample size.

%%%%%%%%%%%%%%%%%%%%%%%%%%%%%%%%%%%%%%%%%%%%%%%%%%%%%%%%%%%%
\subsection{Sampling Variability}
\label{subsec:sampling-variability}
%%%%%%%%%%%%%%%%%%%%%%%%%%%%%%%%%%%%%%%%%%%%%%%%%%%%%%%%%%%%

Different random samples from the same population produce different
statistics.

For example,

\[
\overline{x}_1
\neq
\overline{x}_2
\]

for two different samples.

This natural sample-to-sample variation is called sampling variability.

It is not a mistake.

It is a fundamental source of uncertainty quantified by statistical
inference.

%%%%%%%%%%%%%%%%%%%%%%%%%%%%%%%%%%%%%%%%%%%%%%%%%%%%%%%%%%%%
\subsection{Worked Example: Sampling Variability}
\label{subsec:worked-sampling-variability}
%%%%%%%%%%%%%%%%%%%%%%%%%%%%%%%%%%%%%%%%%%%%%%%%%%%%%%%%%%%%

\begin{workedexamplebox}{Two Samples from the Same Population}

Suppose a population is

\[
\{2,4,6,8\}.
\]

Two samples of size \(2\) might be

\[
\{2,4\}
\]

and

\[
\{6,8\}.
\]

Their means are

\[
3
\]

and

\[
7.
\]

\begin{answerbox}
\[
\boxed{
\text{Different valid samples can produce different statistics.}
}
\]
\end{answerbox}

This is sampling variability.

\end{workedexamplebox}

%%%%%%%%%%%%%%%%%%%%%%%%%%%%%%%%%%%%%%%%%%%%%%%%%%%%%%%%%%%%
\subsection{Sample Size and Precision}
\label{subsec:sample-size-precision}
%%%%%%%%%%%%%%%%%%%%%%%%%%%%%%%%%%%%%%%%%%%%%%%%%%%%%%%%%%%%

Under many standard sampling models, larger samples reduce sampling
variability.

For an independent sample mean,

\[
\operatorname{Var}(\overline{X})
=
\frac{\sigma^2}{n}.
\]

Thus,

\[
\boxed{
\operatorname{SE}(\overline{X})
\propto
\frac1{\sqrt{n}}.
}
\]

However, increasing \(n\) does not automatically reduce systematic
bias.

%%%%%%%%%%%%%%%%%%%%%%%%%%%%%%%%%%%%%%%%%%%%%%%%%%%%%%%%%%%%
\subsection{Margin of Error Preview}
\label{subsec:margin-error-preview}
%%%%%%%%%%%%%%%%%%%%%%%%%%%%%%%%%%%%%%%%%%%%%%%%%%%%%%%%%%%%

Many confidence intervals have the general form

\[
\boxed{
\text{estimate}
\pm
\text{critical value}
\times
\text{standard error}.
}
\]

The second term is often called a margin of error.

Thus larger samples generally produce smaller margins of error because
standard errors shrink.

Confidence intervals are developed in
Section~\ref{sec:confidence-intervals}.

%%%%%%%%%%%%%%%%%%%%%%%%%%%%%%%%%%%%%%%%%%%%%%%%%%%%%%%%%%%%
\subsection{Post-Stratification and Weight Adjustment Preview}
\label{subsec:poststratification-preview}
%%%%%%%%%%%%%%%%%%%%%%%%%%%%%%%%%%%%%%%%%%%%%%%%%%%%%%%%%%%%

Survey weights may be adjusted after data collection so weighted sample
totals align with known population totals.

Examples include adjustments by:

\[
\text{age},
\]

\[
\text{region},
\]

\[
\text{sex},
\]

or

\[
\text{other known categories}.
\]

These procedures can reduce some forms of imbalance but cannot
automatically repair every source of selection bias.

%%%%%%%%%%%%%%%%%%%%%%%%%%%%%%%%%%%%%%%%%%%%%%%%%%%%%%%%%%%%
\subsection{Nonresponse Weighting Preview}
\label{subsec:nonresponse-weighting-preview}
%%%%%%%%%%%%%%%%%%%%%%%%%%%%%%%%%%%%%%%%%%%%%%%%%%%%%%%%%%%%

If response probabilities differ across groups, weights may be adjusted
to compensate.

A schematic adjusted weight is

\[
\boxed{
w_i
=
\frac{
1
}{
\pi_i r_i
},
}
\]

where

\[
\pi_i
\]

is the sampling inclusion probability and

\[
r_i
\]

is an estimated response probability.

Such adjustments depend on assumptions about the nonresponse mechanism.

%%%%%%%%%%%%%%%%%%%%%%%%%%%%%%%%%%%%%%%%%%%%%%%%%%%%%%%%%%%%
\subsection{Selection Diagrams Conceptually}
\label{subsec:selection-diagrams-conceptually}
%%%%%%%%%%%%%%%%%%%%%%%%%%%%%%%%%%%%%%%%%%%%%%%%%%%%%%%%%%%%

One useful way to reason about sampling is to distinguish:

\[
\boxed{
\text{target population},
}
\]

\[
\boxed{
\text{sampling frame},
}
\]

\[
\boxed{
\text{selected sample},
}
\]

and

\[
\boxed{
\text{responding sample}.
}
\]

Bias may enter at every transition.

Thus data quality should be evaluated throughout the full selection
process.

%%%%%%%%%%%%%%%%%%%%%%%%%%%%%%%%%%%%%%%%%%%%%%%%%%%%%%%%%%%%
\subsection{Ethics in Data Collection}
\label{subsec:ethics-data-collection}
%%%%%%%%%%%%%%%%%%%%%%%%%%%%%%%%%%%%%%%%%%%%%%%%%%%%%%%%%%%%

Statistical studies involving people require ethical consideration.

Important principles include:

\[
\boxed{
\text{informed consent},
}
\]

\[
\boxed{
\text{privacy},
}
\]

\[
\boxed{
\text{confidentiality},
}
\]

\[
\boxed{
\text{risk minimization},
}
\]

and

\[
\boxed{
\text{fair participant selection}.
}
\]

Ethical design is part of scientific validity.

%%%%%%%%%%%%%%%%%%%%%%%%%%%%%%%%%%%%%%%%%%%%%%%%%%%%%%%%%%%%
\subsection{Informed Consent}
\label{subsec:informed-consent}
%%%%%%%%%%%%%%%%%%%%%%%%%%%%%%%%%%%%%%%%%%%%%%%%%%%%%%%%%%%%

Participants should generally receive enough information to make an
informed decision about participation.

Relevant information may include:

\[
\text{study purpose},
\]

\[
\text{procedures},
\]

\[
\text{risks},
\]

\[
\text{benefits},
\]

and

\[
\text{data handling}.
\]

Specific legal and institutional requirements vary by context.

%%%%%%%%%%%%%%%%%%%%%%%%%%%%%%%%%%%%%%%%%%%%%%%%%%%%%%%%%%%%
\subsection{Privacy and Confidentiality}
\label{subsec:privacy-confidentiality}
%%%%%%%%%%%%%%%%%%%%%%%%%%%%%%%%%%%%%%%%%%%%%%%%%%%%%%%%%%%%

Privacy concerns whether information is collected or observed.

Confidentiality concerns how collected information is protected and
shared.

Removing obvious identifiers does not always guarantee anonymity.

Combinations of variables may sometimes identify individuals.

%%%%%%%%%%%%%%%%%%%%%%%%%%%%%%%%%%%%%%%%%%%%%%%%%%%%%%%%%%%%
\subsection{Reproducible Data Collection}
\label{subsec:reproducible-data-collection}
%%%%%%%%%%%%%%%%%%%%%%%%%%%%%%%%%%%%%%%%%%%%%%%%%%%%%%%%%%%%

A reproducible study should document:

\begin{itemize}
    \item target population;
    \item eligibility criteria;
    \item sampling frame;
    \item sampling procedure;
    \item recruitment procedure;
    \item measurement definitions;
    \item missing-data rules;
    \item exclusion criteria;
    \item treatment allocation;
    \item data-cleaning procedures.
\end{itemize}

Without this information, statistical results can be difficult to
interpret or reproduce.

%%%%%%%%%%%%%%%%%%%%%%%%%%%%%%%%%%%%%%%%%%%%%%%%%%%%%%%%%%%%
\subsection{Study Design and Validity}
\label{subsec:study-design-validity}
%%%%%%%%%%%%%%%%%%%%%%%%%%%%%%%%%%%%%%%%%%%%%%%%%%%%%%%%%%%%

Two broad forms of validity are often distinguished.

\textbf{Internal validity} concerns whether the study accurately
identifies relationships within the studied setting.

\textbf{External validity} concerns whether conclusions generalize beyond
the observed study units and conditions.

Random assignment primarily strengthens internal causal validity.

Probability sampling primarily strengthens population generalization.

%%%%%%%%%%%%%%%%%%%%%%%%%%%%%%%%%%%%%%%%%%%%%%%%%%%%%%%%%%%%
\subsection{Internal Validity}
\label{subsec:internal-validity}
%%%%%%%%%%%%%%%%%%%%%%%%%%%%%%%%%%%%%%%%%%%%%%%%%%%%%%%%%%%%

Threats to internal validity include:

\[
\boxed{
\text{confounding},
}
\]

\[
\boxed{
\text{differential attrition},
}
\]

\[
\boxed{
\text{measurement changes},
}
\]

and

\[
\boxed{
\text{treatment contamination}.
}
\]

A randomized experiment is designed to reduce several of these threats.

%%%%%%%%%%%%%%%%%%%%%%%%%%%%%%%%%%%%%%%%%%%%%%%%%%%%%%%%%%%%
\subsection{External Validity}
\label{subsec:external-validity}
%%%%%%%%%%%%%%%%%%%%%%%%%%%%%%%%%%%%%%%%%%%%%%%%%%%%%%%%%%%%

External validity depends on how well the study conditions reflect the
settings to which conclusions are applied.

Relevant issues include:

\[
\boxed{
\text{population selection},
}
\]

\[
\boxed{
\text{geographic setting},
}
\]

\[
\boxed{
\text{time period},
}
\]

and

\[
\boxed{
\text{treatment implementation}.
}
\]

A perfectly randomized experiment can still have limited external
validity.

%%%%%%%%%%%%%%%%%%%%%%%%%%%%%%%%%%%%%%%%%%%%%%%%%%%%%%%%%%%%
\subsection{Causal Conclusions}
\label{subsec:causal-conclusions-data-collection}
%%%%%%%%%%%%%%%%%%%%%%%%%%%%%%%%%%%%%%%%%%%%%%%%%%%%%%%%%%%%

Causal inference requires more than observing an association.

A randomized controlled experiment provides strong causal evidence
because treatment assignment is randomized.

In observational studies, causal interpretation generally requires
additional assumptions, design strategies, or statistical adjustments.

Thus:

\[
\boxed{
\text{association}
\not\Rightarrow
\text{causation}.
}
\]

%%%%%%%%%%%%%%%%%%%%%%%%%%%%%%%%%%%%%%%%%%%%%%%%%%%%%%%%%%%%
\subsection{Potential Outcomes Preview}
\label{subsec:potential-outcomes-preview}
%%%%%%%%%%%%%%%%%%%%%%%%%%%%%%%%%%%%%%%%%%%%%%%%%%%%%%%%%%%%

A useful formal framework for causal inference associates each unit with
potential responses under different treatments.

For two treatments,

\[
Y_i(1)
\]

may denote the response unit \(i\) would have under treatment, while

\[
Y_i(0)
\]

denotes the response under control.

The individual causal effect is conceptually

\[
\boxed{
Y_i(1)-Y_i(0).
}
\]

Only one potential outcome can usually be observed for the same unit at
the same time.

This is the fundamental missing-data problem of causal inference.

%%%%%%%%%%%%%%%%%%%%%%%%%%%%%%%%%%%%%%%%%%%%%%%%%%%%%%%%%%%%
\subsection{Average Treatment Effect Preview}
\label{subsec:average-treatment-effect-preview}
%%%%%%%%%%%%%%%%%%%%%%%%%%%%%%%%%%%%%%%%%%%%%%%%%%%%%%%%%%%%

The average treatment effect is conceptually

\[
\boxed{
ATE
=
\mathbb{E}
[
Y(1)-Y(0)
].
}
\]

Random assignment helps make the treatment and control groups comparable
so differences in average observed outcomes can estimate this causal
quantity.

A full treatment belongs to causal inference and experimental design.

%%%%%%%%%%%%%%%%%%%%%%%%%%%%%%%%%%%%%%%%%%%%%%%%%%%%%%%%%%%%
\subsection{Data Collection Workflow}
\label{subsec:data-collection-workflow}
%%%%%%%%%%%%%%%%%%%%%%%%%%%%%%%%%%%%%%%%%%%%%%%%%%%%%%%%%%%%

A systematic workflow is:

\begin{enumerate}
    \item define the research question;
    \item define the target population;
    \item identify the observational units;
    \item identify the variables to measure;
    \item decide whether an observational study or experiment is
          appropriate;
    \item construct or evaluate the sampling frame;
    \item choose a sampling design;
    \item determine sample size;
    \item define measurement procedures;
    \item plan treatment assignment if conducting an experiment;
    \item anticipate nonresponse and missing data;
    \item document inclusion and exclusion criteria;
    \item conduct quality-control checks;
    \item preserve information about weights and design variables;
    \item assess limitations before making inferential claims.
\end{enumerate}

%%%%%%%%%%%%%%%%%%%%%%%%%%%%%%%%%%%%%%%%%%%%%%%%%%%%%%%%%%%%
\subsection{Choosing a Sampling Design}
\label{subsec:choosing-sampling-design}
%%%%%%%%%%%%%%%%%%%%%%%%%%%%%%%%%%%%%%%%%%%%%%%%%%%%%%%%%%%%

A simple random sample is attractive when:

\[
\boxed{
\text{a complete sampling frame is available}
}
\]

and sampling individual units is practical.

Stratified sampling is attractive when:

\[
\boxed{
\text{important subgroups must be represented}.
}
\]

Cluster sampling is attractive when:

\[
\boxed{
\text{population units are geographically or organizationally grouped}.
}
\]

Multistage sampling is attractive for large, complex populations.

%%%%%%%%%%%%%%%%%%%%%%%%%%%%%%%%%%%%%%%%%%%%%%%%%%%%%%%%%%%%
\subsection{Common Errors}
\label{subsec:data-collection-common-errors}
%%%%%%%%%%%%%%%%%%%%%%%%%%%%%%%%%%%%%%%%%%%%%%%%%%%%%%%%%%%%

\subsubsection{Assuming a Large Sample Must Be Representative}

A very large convenience sample can remain badly biased.

Sample size primarily affects random error, not systematic selection
bias.

%%%%%%%%%%%%%%%%%%%%%%%%%%%%%%%%%%%%%%%%%%%%%%%%%%%%%%%%%%%%
\subsubsection{Confusing Random Sampling with Random Assignment}

Random sampling supports generalization.

Random assignment supports causal comparison.

They solve different statistical problems.

%%%%%%%%%%%%%%%%%%%%%%%%%%%%%%%%%%%%%%%%%%%%%%%%%%%%%%%%%%%%
\subsubsection{Calling a Sample Random Because Participation Was Unpredictable}

A sample is not a probability sample merely because the investigator did
not know who would participate.

The selection mechanism itself must contain a defined randomization
procedure.

%%%%%%%%%%%%%%%%%%%%%%%%%%%%%%%%%%%%%%%%%%%%%%%%%%%%%%%%%%%%
\subsubsection{Assuming Every Population Unit Must Have Equal Selection Probability}

Probability sampling requires known selection probabilities, but they
need not always be equal.

Unequal-probability designs can be valid when properly weighted.

%%%%%%%%%%%%%%%%%%%%%%%%%%%%%%%%%%%%%%%%%%%%%%%%%%%%%%%%%%%%
\subsubsection{Confusing Strata and Clusters}

Stratified designs sample from every stratum.

Cluster designs sample only selected clusters.

%%%%%%%%%%%%%%%%%%%%%%%%%%%%%%%%%%%%%%%%%%%%%%%%%%%%%%%%%%%%
\subsubsection{Ignoring Sampling Weights}

If units have unequal inclusion probabilities, unweighted estimates may
not represent the population correctly.

%%%%%%%%%%%%%%%%%%%%%%%%%%%%%%%%%%%%%%%%%%%%%%%%%%%%%%%%%%%%
\subsubsection{Treating Nonresponse as Random Automatically}

Nonresponse may be related to the outcome or important covariates.

Ignoring this mechanism can create bias.

%%%%%%%%%%%%%%%%%%%%%%%%%%%%%%%%%%%%%%%%%%%%%%%%%%%%%%%%%%%%
\subsubsection{Assuming a Census Has No Error}

A census may still contain coverage, measurement, processing, and
nonresponse errors.

%%%%%%%%%%%%%%%%%%%%%%%%%%%%%%%%%%%%%%%%%%%%%%%%%%%%%%%%%%%%
\subsubsection{Removing Nonrespondents from the Target Population}

Failure to respond does not usually change who belongs to the target
population.

It changes which target units are observed.

%%%%%%%%%%%%%%%%%%%%%%%%%%%%%%%%%%%%%%%%%%%%%%%%%%%%%%%%%%%%
\subsubsection{Confusing Confounding with Sampling Bias}

Sampling bias concerns how units enter the sample.

Confounding concerns the inability to separate effects of explanatory
variables.

They are distinct problems.

%%%%%%%%%%%%%%%%%%%%%%%%%%%%%%%%%%%%%%%%%%%%%%%%%%%%%%%%%%%%
\subsubsection{Assuming Random Assignment Creates Identical Groups}

Randomization balances characteristics probabilistically, not perfectly
in every realized experiment.

%%%%%%%%%%%%%%%%%%%%%%%%%%%%%%%%%%%%%%%%%%%%%%%%%%%%%%%%%%%%
\subsubsection{Ignoring Clustering in Analysis}

Observations from the same cluster may be dependent.

Treating them as independent can underestimate uncertainty.

%%%%%%%%%%%%%%%%%%%%%%%%%%%%%%%%%%%%%%%%%%%%%%%%%%%%%%%%%%%%
\subsubsection{Ignoring the Original Study Design}

Statistical analysis should preserve information about:

\[
\text{strata},
\]

\[
\text{clusters},
\]

\[
\text{weights},
\]

and

\[
\text{selection probabilities}.
\]

Ignoring these may produce incorrect standard errors or estimates.

%%%%%%%%%%%%%%%%%%%%%%%%%%%%%%%%%%%%%%%%%%%%%%%%%%%%%%%%%%%%
\subsection{Applications}
\label{subsec:data-collection-applications}
%%%%%%%%%%%%%%%%%%%%%%%%%%%%%%%%%%%%%%%%%%%%%%%%%%%%%%%%%%%%

\begin{applicationbox}

\textbf{Public Opinion Surveys.}
Probability sampling, weighting, and nonresponse adjustment are used to
estimate population opinions.

\medskip

\textbf{Government Statistics.}
Household, labor, health, and economic surveys commonly use stratified
and multistage probability designs.

\medskip

\textbf{Education Research.}
Students may be sampled within classrooms or schools, creating clustered
data structures.

\medskip

\textbf{Clinical Research.}
Randomized experiments compare treatments while controlling potential
confounding through assignment.

\medskip

\textbf{Environmental Science.}
Sampling locations may be stratified geographically or selected through
spatial designs.

\medskip

\textbf{Manufacturing.}
Quality-control samples are taken from production processes to estimate
defect rates and process characteristics.

\medskip

\textbf{Epidemiology.}
Cohort, case-control, cross-sectional, and randomized studies address
different scientific questions.

\medskip

\textbf{Data Science.}
Large observational datasets require careful evaluation of selection,
measurement, missingness, and representativeness before modeling.

\end{applicationbox}

%%%%%%%%%%%%%%%%%%%%%%%%%%%%%%%%%%%%%%%%%%%%%%%%%%%%%%%%%%%%
\subsection{Exercises}
\label{subsec:data-collection-exercises}
%%%%%%%%%%%%%%%%%%%%%%%%%%%%%%%%%%%%%%%%%%%%%%%%%%%%%%%%%%%%

\begin{exercise}[1]
Define the target population.
\end{exercise}

\begin{exercise}[1]
Define a sampling frame.
\end{exercise}

\begin{exercise}[1]
Define a probability sample.
\end{exercise}

\begin{exercise}[1]
Define a simple random sample.
\end{exercise}

\begin{exercise}[1]
Distinguish sampling with replacement from sampling without replacement.
\end{exercise}

\begin{exercise}[1]
Define stratified sampling.
\end{exercise}

\begin{exercise}[1]
Define cluster sampling.
\end{exercise}

\begin{exercise}[1]
Define systematic sampling.
\end{exercise}

\begin{exercise}[1]
Define sampling error.
\end{exercise}

\begin{exercise}[1]
Define nonsampling error.
\end{exercise}

\begin{exercise}[1]
Define undercoverage.
\end{exercise}

\begin{exercise}[1]
Define nonresponse bias.
\end{exercise}

\begin{exercise}[1]
Define an observational study.
\end{exercise}

\begin{exercise}[1]
Define an experiment.
\end{exercise}

\begin{exercise}[1]
Distinguish random sampling from random assignment.
\end{exercise}

\begin{exercise}[1]
Define confounding.
\end{exercise}

\begin{exercise}[1]
Define replication, randomization, and blocking.
\end{exercise}

\begin{exercise}[2]
A population contains \(20\) units. How many distinct simple random
samples of size \(4\) are possible?
\end{exercise}

\begin{exercise}[2]
For an SRS of size \(10\) from a population of size \(100\), find the
inclusion probability of a particular unit.
\end{exercise}

\begin{exercise}[2]
A population has \(N=500\) units and a sample has \(n=50\). Compute the
sampling fraction.
\end{exercise}

\begin{exercise}[2]
Compute the finite-population correction

\[
\sqrt{
\frac{N-n}{N-1}
}
\]

for

\[
N=100,
\qquad
n=20.
\]
\end{exercise}

\begin{exercise}[2]
A population contains three strata of sizes

\[
400,\ 350,\ 250.
\]

Construct a proportional stratified sample of total size \(100\).
\end{exercise}

\begin{exercise}[2]
Explain why oversampling a small subgroup may be useful.
\end{exercise}

\begin{exercise}[2]
A systematic sample selects every \(10\)-th record after a random start.
Explain one situation in which this design could be problematic.
\end{exercise}

\begin{exercise}[2]
Classify each as a likely example of undercoverage, nonresponse, or
response bias:
\begin{itemize}
    \item selected participants refuse to answer;
    \item a survey frame omits people without internet access;
    \item participants underreport a socially sensitive behavior.
\end{itemize}
\end{exercise}

\begin{exercise}[2]
For an SRS, explain why every population unit has inclusion probability

\[
n/N.
\]
\end{exercise}

\begin{exercise}[2]
Suppose a unit has inclusion probability

\[
\pi_i=0.05.
\]

Find its basic sampling weight.
\end{exercise}

\begin{exercise}[2]
Identify the experimental units, explanatory variable, and response
variable in an experiment comparing two fertilizers on crop yield.
\end{exercise}

\begin{exercise}[2]
Explain how blocking by soil type could improve the fertilizer
experiment.
\end{exercise}

\begin{exercise}[2]
Give an example of a matched-pairs experiment.
\end{exercise}

\begin{exercise}[2]
Distinguish a prospective study from a retrospective study.
\end{exercise}

\begin{exercise}[2]
Explain why a large online voluntary-response poll need not represent the
general population.
\end{exercise}

\begin{exercise}[3]
Prove that each unit in an SRS of size \(n\) from \(N\) units has
inclusion probability

\[
n/N.
\]
\end{exercise}

\begin{exercise}[3]
Derive the probability that two specified population units are both
included in an SRS of size \(n\).
\end{exercise}

\begin{exercise}[3]
Show that sampling indicators in an SRS without replacement are
negatively correlated.
\end{exercise}

\begin{exercise}[3]
Derive the finite-population correction for the variance of an SRS sample
mean.
\end{exercise}

\begin{exercise}[3]
Show that the Horvitz--Thompson estimator is unbiased for the population
total.
\end{exercise}

\begin{exercise}[3]
Derive a stratified estimator of the population mean:

\[
\overline{y}_{st}
=
\sum_{h=1}^{H}
\frac{N_h}{N}
\overline{y}_h.
\]
\end{exercise}

\begin{exercise}[3]
Explain why stratification can reduce variance when within-stratum
variation is small.
\end{exercise}

\begin{exercise}[3]
Explain why cluster sampling may increase variance when within-cluster
correlation is positive.
\end{exercise}

\begin{exercise}[3]
For a two-stage design, write the overall inclusion probability as a
product of stage-specific probabilities.
\end{exercise}

\begin{exercise}[3]
Construct an example where random sampling occurs but random assignment
does not.
\end{exercise}

\begin{exercise}[3]
Construct an example where random assignment occurs but random sampling
does not.
\end{exercise}

\begin{exercise}[3]
Explain what conclusions each of the previous two designs can support.
\end{exercise}

\begin{exercise}[3]
Construct a causal diagram conceptually involving treatment, outcome,
and a confounder.
\end{exercise}

\begin{exercise}[3]
Explain why adjusting for sample size alone cannot remove selection
bias.
\end{exercise}

\begin{exercise}[3]
Compare internal and external validity.
\end{exercise}

\begin{exercise}[3]
Explain why repeated measurements on the same individual should not
usually be treated as independent observations.
\end{exercise}

\begin{exercise}[3]
Explain why case-control data do not directly estimate population disease
prevalence from the fraction of sampled cases.
\end{exercise}

\begin{exercise}[3]
Compare a census with a well-designed probability sample in terms of
cost, sampling error, and nonsampling error.
\end{exercise}

\begin{exercise}[3]
Explain how a weighted estimator can correct unequal inclusion
probabilities.
\end{exercise}

\begin{exercise}[3]
Explain why weighting cannot necessarily remove bias caused by unmeasured
differences between respondents and nonrespondents.
\end{exercise}

\begin{exercise}[4]
Study unequal-probability sampling and derive Horvitz--Thompson variance
estimators.
\end{exercise}

\begin{exercise}[4]
Study Neyman allocation for stratified sampling.
\end{exercise}

\begin{exercise}[4]
Investigate probability proportional to size sampling.
\end{exercise}

\begin{exercise}[4]
Study multistage cluster sampling and design effects.
\end{exercise}

\begin{exercise}[4]
Investigate intraclass correlation and effective sample size.
\end{exercise}

\begin{exercise}[4]
Study calibration weighting and post-stratification.
\end{exercise}

\begin{exercise}[4]
Investigate inverse-probability weighting for nonresponse.
\end{exercise}

\begin{exercise}[4]
Study missing completely at random, missing at random, and missing not at
random mechanisms.
\end{exercise}

\begin{exercise}[4]
Investigate the potential-outcomes framework for causal inference.
\end{exercise}

\begin{exercise}[4]
Study randomized-block and factorial experiments.
\end{exercise}

\begin{exercise}[4]
Investigate cluster-randomized trials.
\end{exercise}

\begin{exercise}[4]
Study observational-study adjustment using matching and propensity
scores.
\end{exercise}

\begin{exercise}[4]
Investigate survey nonresponse and modern weighting methods.
\end{exercise}

\begin{exercise}[4]
Study the distinction between design-based, model-based, and
model-assisted survey inference.
\end{exercise}

%%%%%%%%%%%%%%%%%%%%%%%%%%%%%%%%%%%%%%%%%%%%%%%%%%%%%%%%%%%%
\subsection{Conceptual Summary}
\label{subsec:data-collection-summary}
%%%%%%%%%%%%%%%%%%%%%%%%%%%%%%%%%%%%%%%%%%%%%%%%%%%%%%%%%%%%

A target population is the population to which conclusions are intended
to apply.

A sampling frame is the mechanism used to identify possible sample units.

A probability sample uses a known random selection mechanism.

In a simple random sample of size \(n\) from \(N\) units,

\[
\boxed{
P(\text{unit }i\text{ selected})
=
\frac nN.
}
\]

The sampling fraction is

\[
\boxed{
f=\frac nN.
}
\]

Sampling without replacement may require the finite-population
correction

\[
\boxed{
\sqrt{
\frac{N-n}{N-1}
}.
}
\]

Stratified sampling divides the population into strata and samples from
each.

Cluster sampling selects groups of units.

Multistage sampling selects units through several randomized levels.

If inclusion probabilities differ,

\[
\boxed{
w_i
=
\frac1{\pi_i}
}
\]

provides the basic design weight.

Sampling error is caused by observing a sample instead of the entire
population.

Nonsampling error includes

\[
\boxed{
\text{coverage}
+
\text{nonresponse}
+
\text{measurement}
+
\text{processing errors}.
}
\]

Random sampling primarily supports

\[
\boxed{
\text{population generalization}.
}
\]

Random assignment primarily supports

\[
\boxed{
\text{causal comparison}.
}
\]

A strong statistical study therefore depends on both mathematical
analysis and the design that created the data.

%%%%%%%%%%%%%%%%%%%%%%%%%%%%%%%%%%%%%%%%%%%%%%%%%%%%%%%%%%%%
\subsection{Section Connections}
\label{subsec:data-collection-connections}
%%%%%%%%%%%%%%%%%%%%%%%%%%%%%%%%%%%%%%%%%%%%%%%%%%%%%%%%%%%%

\chapterconnection{
Data Collection and Sampling determines the mechanism by which the
descriptive summaries of Section~\ref{sec:descriptive-statistics} become
available for analysis. Probability sampling creates random statistics,
and their sample-to-sample variation leads directly to the sampling
distributions developed in the next section. Random assignment provides
the design foundation for causal comparisons and is developed more
fully in Section~\ref{sec:experimental-design}. Sampling weights,
nonresponse, clustering, and stratification remain important throughout
estimation, confidence intervals, hypothesis testing, regression, and
survey analysis. The next section, Sampling Distributions, formalizes
the probability distributions of statistics generated by repeated
sampling.
}

%%%%%%%%%%%%%%%%%%%%%%%%%%%%%%%%%%%%%%%%%%%%%%%%%%%%%%%%%%%%
% SECTION SOURCE TRACKING
%%%%%%%%%%%%%%%%%%%%%%%%%%%%%%%%%%%%%%%%%%%%%%%%%%%%%%%%%%%%

%------------------------------------------------------------
% 8.2 Data Collection and Sampling
%------------------------------------------------------------
%
% Source basis:
%   - Standard survey-sampling theory.
%   - Standard introductory statistical study design.
%   - Standard experimental-design and observational-study concepts.
%
% Used for:
%   - target and study populations;
%   - sampling frames;
%   - sampling and observational units;
%   - census and sample surveys;
%   - probability and nonprobability sampling;
%   - convenience and voluntary-response samples;
%   - simple random sampling;
%   - sampling with and without replacement;
%   - inclusion probabilities;
%   - finite-population corrections;
%   - systematic sampling;
%   - stratified sampling;
%   - proportional and disproportionate allocation;
%   - Neyman allocation preview;
%   - cluster sampling;
%   - multistage sampling;
%   - cluster correlation;
%   - design effects;
%   - unequal-probability sampling;
%   - sampling weights;
%   - Horvitz--Thompson estimator preview;
%   - sampling and nonsampling error;
%   - selection bias;
%   - undercoverage;
%   - nonresponse;
%   - response bias;
%   - measurement error;
%   - validity and reliability;
%   - observational studies;
%   - experiments;
%   - experimental units;
%   - treatments;
%   - response and explanatory variables;
%   - confounding;
%   - random assignment;
%   - control groups and placebo;
%   - blinding;
%   - replication;
%   - blocking;
%   - matched pairs;
%   - cross-sectional and longitudinal studies;
%   - prospective and retrospective studies;
%   - case-control and cohort studies;
%   - administrative and secondary data;
%   - design-based and model-based inference;
%   - sampling variability;
%   - standard-error preview;
%   - weighting adjustments;
%   - internal and external validity;
%   - causal-inference preview;
%   - potential-outcomes preview;
%   - ethics and reproducibility;
%   - applications and exercises.
%
% Notes:
%   - Sampling Distributions are developed next in Section 8.3.
%   - Experimental Design is developed more fully in Section 8.9.
%   - Confidence intervals and hypothesis tests later depend on the
%     sampling mechanisms introduced here.
%   - Advanced survey-sampling estimators are introduced only as previews
%     in this section.
%
%%%%%%%%%%%%%%%%%%%%%%%%%%%%%%%%%%%%%%%%%%%%%%%%%%%%%%%%%%%%